% GENERATED FILE. Edit canonical lessons, then run npm run generate:books.
% canonical-source-hash: fnv1a64:8f514f128dd50a9a
% canonical-lessons: PA-C229-gandh, PA-C229-bannhna, PA-C229-kassna, PA-C229-dori, PA-C229-ghundi

\chapter{Knots and Strings}
\label{ch:pa-a2-knots-and-strings}
\hlchaptermodality{229}

\begin{chapteropening}
\textbf{By the end of this chapter:} \emph{I can say what I tie and tighten.}

Complete the last lesson of chapter 229: I can say what I tie and tighten.
\end{chapteropening}

\section[\texorpdfstring{\pa{ਗੰਢ} (gaṇḍh)}{ਗੰਢ (gaṇḍh)}]{\pa{ਗੰਢ} (gaṇḍh) — a knot}

\label{lesson:PA-C229-gandh}



\begin{quote}
Before the new one: say the Punjabi for a rope, then the Punjabi for to wipe.
\end{quote}

\subsection*{You'll want to know: \pa{ਗੰਢ}}
\textbf{\pa{ਗੰਢ}} — \emph{gaṇḍh} — \textquotedblleft{}a knot\textquotedblright{}.

\begin{quote}
\pa{ਰੱਸੀ} \pa{ਵਿੱਚ} \pa{ਇੱਕ} \pa{ਗੰਢ} \pa{ਹੈ।}
\end{quote}

\begin{quote}
\emph{rassī vich ikk gaṇḍh hai.}
\end{quote}

\begin{quote}
There is a knot in the rope.
\end{quote}

A feminine noun.

\begin{cousinweb}
\textbf{\pa{ਗੰਢ}} (\emph{gaṇḍh}) is inherited from Sanskrit \emph{granthi}, \textquotedblleft{}knot, bunch\textquotedblright{}; CDIAL lists Punjabi \emph{gaṇḍh}, \textquotedblleft{}knot, joint\textquotedblright{} (CDIAL 4354). CDIAL marks the Punjabi row masculine and the Lahnda row feminine; here \pa{ਗੰਢ} is taught as feminine, as the Lahnda row has it.
\end{cousinweb}

\begin{grammarlens}[title={What you've built}]
One more word for saying what you do.
\end{grammarlens}

\subsection*{Your turn}
\begin{itemize}
\raggedright
  \item \emph{Say it:} \emph{gaṇḍh}
  \item \emph{Say it:} \emph{gaṇḍh}, once more
  \item \emph{Say it:} \emph{rassī vich ikk gaṇḍh hai.} — There is a knot in the rope.
  \item \emph{Recall:} say the Punjabi for a rope, then the Punjabi for to wipe, then say \emph{gaṇḍh} again
\end{itemize}

\begin{tcolorbox}[breakable,colback=teal!4,colframe=teal!35!black,title={Before you move on}]
What does \pa{ਗੰਢ} mean? (\textquotedblleft{}A knot\textquotedblright{}.) Say it once more.
\end{tcolorbox}
\section[\texorpdfstring{\pa{ਬੰਨ੍ਹਣਾ} (bannhṇā)}{ਬੰਨ੍ਹਣਾ (bannhṇā)}]{\pa{ਬੰਨ੍ਹਣਾ} (bannhṇā) — to tie}

\label{lesson:PA-C229-bannhna}



\begin{quote}
Before the new one: say the Punjabi for to wipe, then the Punjabi for a knot.
\end{quote}

\subsection*{You'll want to know: \pa{ਬੰਨ੍ਹਣਾ}}
\textbf{\pa{ਬੰਨ੍ਹਣਾ}} — \emph{bannhṇā} — \textquotedblleft{}to tie\textquotedblright{}.

\begin{quote}
\pa{ਮੈਂ} \pa{ਰੱਸੀ} \pa{ਨਾਲ} \pa{ਡੱਬਾ} \pa{ਬੰਨ੍ਹਦੀ} \pa{ਹਾਂ।}
\end{quote}

\begin{quote}
\emph{maiṁ rassī nāl ḍabbā bannhdī hāṁ.}
\end{quote}

\begin{quote}
I tie up the box with a rope.
\end{quote}

Past (I …): \pa{ਮੈਂ} \pa{ਗੰਢ} \pa{ਬੰਨ੍ਹੀ}, \emph{maiṁ gaṇḍh bannhī}.

\begin{cousinweb}
\textbf{\pa{ਬੰਨ੍ਹਣਾ}} (\emph{bannhṇā}) is inherited from Sanskrit \emph{bandhati}, \textquotedblleft{}binds\textquotedblright{}; CDIAL lists Punjabi \emph{bannhṇā}, \textquotedblleft{}to tie, bind\textquotedblright{} (CDIAL 9139).
\end{cousinweb}

\begin{grammarlens}[title={What you've built}]
One more word for saying what you do.
\end{grammarlens}

\subsection*{Your turn}
\begin{itemize}
\raggedright
  \item \emph{Say it:} \emph{bannhṇā}
  \item \emph{Say it:} \emph{bannhṇā}, once more
  \item \emph{Say it:} \emph{maiṁ rassī nāl ḍabbā bannhdī hāṁ.} — I tie up the box with a rope.
  \item \emph{Recall:} say the Punjabi for to wipe, then the Punjabi for a knot, then say \emph{bannhṇā} again
  \item \emph{Write it:} copy \textbf{\pa{ਬੰਨ੍ਹਣਾ}} once, checking the whole word
\end{itemize}

\begin{tcolorbox}[breakable,colback=teal!4,colframe=teal!35!black,title={Before you move on}]
What does \pa{ਬੰਨ੍ਹਣਾ} mean? (\textquotedblleft{}To tie\textquotedblright{}.) Say it once more.
\end{tcolorbox}
\section[\texorpdfstring{\pa{ਕੱਸਣਾ} (kassṇā)}{ਕੱਸਣਾ (kassṇā)}]{\pa{ਕੱਸਣਾ} (kassṇā) — to tighten}

\label{lesson:PA-C229-kassna}



\begin{quote}
Before the new one: say the Punjabi for a knot, then the Punjabi for to tie.
\end{quote}

\subsection*{You'll want to know: \pa{ਕੱਸਣਾ}}
\textbf{\pa{ਕੱਸਣਾ}} — \emph{kassṇā} — \textquotedblleft{}to tighten\textquotedblright{}.

\begin{quote}
\pa{ਮੈਂ} \pa{ਸਵੇਰੇ} \pa{ਪੇਟੀ} \pa{ਕੱਸਦਾ} \pa{ਹਾਂ।}
\end{quote}

\begin{quote}
\emph{maiṁ savere peṭī kassdā hāṁ.}
\end{quote}

\begin{quote}
In the morning I tighten my belt.
\end{quote}

Past (I …): \pa{ਮੈਂ} \pa{ਗੰਢ} \pa{ਕੱਸੀ}, \emph{maiṁ gaṇḍh kassī}.

\begin{cousinweb}
\textbf{\pa{ਕੱਸਣਾ}} (\emph{kassṇā}) is inherited from Sanskrit \emph{karṣati}, \textquotedblleft{}draws, pulls\textquotedblright{}; CDIAL lists Punjabi \emph{kassṇā}, \textquotedblleft{}to tighten\textquotedblright{} (CDIAL 2908).
\end{cousinweb}

\begin{grammarlens}[title={What you've built}]
One more word for saying what you do.
\end{grammarlens}

\subsection*{Your turn}
\begin{itemize}
\raggedright
  \item \emph{Say it:} \emph{kassṇā}
  \item \emph{Say it:} \emph{kassṇā}, once more
  \item \emph{Say it:} \emph{maiṁ savere peṭī kassdā hāṁ.} — In the morning I tighten my belt.
  \item \emph{Recall:} say the Punjabi for a knot, then the Punjabi for to tie, then say \emph{kassṇā} again
  \item \emph{Write it:} copy \textbf{\pa{ਕੱਸਣਾ}} once, checking the whole word
\end{itemize}

\begin{tcolorbox}[breakable,colback=teal!4,colframe=teal!35!black,title={Before you move on}]
What does \pa{ਕੱਸਣਾ} mean? (\textquotedblleft{}To tighten\textquotedblright{}.) Say it once more.
\end{tcolorbox}
\section[\texorpdfstring{\pa{ਡੋਰੀ} (ḍorī)}{ਡੋਰੀ (ḍorī)}]{\pa{ਡੋਰੀ} (ḍorī) — a cord; a drawstring}

\label{lesson:PA-C229-dori}



\begin{quote}
Before the new one: say the Punjabi for to tie, then the Punjabi for to tighten.
\end{quote}

\subsection*{You'll want to know: \pa{ਡੋਰੀ}}
\textbf{\pa{ਡੋਰੀ}} — \emph{ḍorī} — \textquotedblleft{}a cord; a drawstring\textquotedblright{}.

\begin{quote}
\pa{ਸਲਵਾਰ} \pa{ਦੀ} \pa{ਡੋਰੀ} \pa{ਲੰਮੀ} \pa{ਹੈ।}
\end{quote}

\begin{quote}
\emph{salvār dī ḍorī lammī hai.}
\end{quote}

\begin{quote}
The drawstring of the salwar is long.
\end{quote}

A feminine noun.

\begin{cousinweb}
\textbf{\pa{ਡੋਰੀ}} (\emph{ḍorī}) is inherited from Sanskrit \emph{davara}, \textquotedblleft{}string\textquotedblright{}; CDIAL lists Punjabi \emph{ḍorī}, \textquotedblleft{}string\textquotedblright{} (CDIAL 6225).
\end{cousinweb}

\begin{grammarlens}[title={What you've built}]
One more word for saying what you do.
\end{grammarlens}

\subsection*{Your turn}
\begin{itemize}
\raggedright
  \item \emph{Say it:} \emph{ḍorī}
  \item \emph{Say it:} \emph{ḍorī}, once more
  \item \emph{Say it:} \emph{salvār dī ḍorī lammī hai.} — The drawstring of the salwar is long.
  \item \emph{Recall:} say the Punjabi for to tie, then the Punjabi for to tighten, then say \emph{ḍorī} again
  \item \emph{Write it:} copy \textbf{\pa{ਡੋਰੀ}} once, checking the whole word
\end{itemize}

\begin{tcolorbox}[breakable,colback=teal!4,colframe=teal!35!black,title={Before you move on}]
What does \pa{ਡੋਰੀ} mean? (\textquotedblleft{}A cord; a drawstring\textquotedblright{}.) Say it once more.
\end{tcolorbox}
\section[\texorpdfstring{\pa{ਘੁੰਡੀ} (ghuṇḍī)}{ਘੁੰਡੀ (ghuṇḍī)}]{\pa{ਘੁੰਡੀ} (ghuṇḍī) — a knotted cloth button}

\label{lesson:PA-C229-ghundi}



\begin{quote}
Before the new one: say the Punjabi for to tighten, then the Punjabi for a cord.
\end{quote}

\subsection*{You'll want to know: \pa{ਘੁੰਡੀ}}
\textbf{\pa{ਘੁੰਡੀ}} — \emph{ghuṇḍī} — \textquotedblleft{}a knotted cloth button\textquotedblright{}.

\begin{quote}
\pa{ਮੇਰੇ} \pa{ਕੁੜਤੇ} \pa{ਉੱਤੇ} \pa{ਤਿੰਨ} \pa{ਘੁੰਡੀਆਂ} \pa{ਹਨ।}
\end{quote}

\begin{quote}
\emph{mere kuṛte utte tinn ghuṇḍīāṁ han.}
\end{quote}

\begin{quote}
My kurta has three cloth buttons.
\end{quote}

A feminine noun.

\begin{cousinweb}
\textbf{\pa{ਘੁੰਡੀ}} (\emph{ghuṇḍī}) is inherited from a reconstructed \emph{*ghuṇṭa}, \textquotedblleft{}knot, tag, button\textquotedblright{}; CDIAL lists Punjabi \emph{ghuṇḍī}, \textquotedblleft{}knot, button\textquotedblright{} (CDIAL 4483).
\end{cousinweb}

\begin{grammarlens}[title={What you've built}]
One more word for saying what you do.
\end{grammarlens}

\subsection*{Your turn}
\begin{itemize}
\raggedright
  \item \emph{Say it:} \emph{ghuṇḍī}
  \item \emph{Say it:} \emph{ghuṇḍī}, once more
  \item \emph{Say it:} \emph{mere kuṛte utte tinn ghuṇḍīāṁ han.} — My kurta has three cloth buttons.
  \item \emph{Recall:} say the Punjabi for to tighten, then the Punjabi for a cord, then say \emph{ghuṇḍī} again
  \item \emph{Write it:} copy \textbf{\pa{ਘੁੰਡੀ}} once, checking the whole word
\end{itemize}

\begin{tcolorbox}[breakable,colback=teal!4,colframe=teal!35!black,title={Before you move on}]
What does \pa{ਘੁੰਡੀ} mean? (\textquotedblleft{}A knotted cloth button\textquotedblright{}.) Say it once more.
\end{tcolorbox}
